\chapter{The Fine Structure of a Brownian Path}

\begin{center}\itshape
Zoom in on a Brownian path and it looks the same: no tangent, no smooth piece, the same jagged lace at every scale.\\
But the zoom has an exact law. Increments over a time $h$ are never larger than $\sqrt{2h\log(1/h)}$, and the times\\
when the path returns to zero form a dust of dimension exactly one half.
\end{center}

Paul L\'evy did not compute Brownian motion; he looked at it. Between 1937 and 1948, in two books that his contemporaries found hard to read and impossible to
put down, he described the path as if he had held it in his hands: how fast it can oscillate, how its zeros cluster, how its maximum and its current position are
related, how long it spends on each side of zero. Michel Lo\`eve, who knew him well, wrote in his obituary of L\'evy that he worked by an intuition so strong that proofs
came afterwards, sometimes much afterwards, and sometimes from others (we paraphrase). Many of L\'evy's statements were proved rigorously only after the strong Markov
property had been made precise in the 1950s, and several of them are the subject of this chapter.

L\'evy's most famous discovery about the fine structure is the exact size of the oscillations. For a single fixed time $t$, the law of the iterated logarithm of Volume~11
says that $|B_{t+h}-B_t|$ is at most about $\sqrt{2h\log\log(1/h)}$ as $h\to0$. But there are uncountably many times, and some of them must be exceptional. L\'evy found in
1937 that the worst case over all times replaces $\log\log$ by $\log$: the uniform modulus of continuity is $\sqrt{2h\log(1/h)}$, with the constant exactly $\sqrt2$. The path is
H\"older of every order below $\frac12$, as Chapter~1 showed, and the logarithm measures precisely how far it is from order $\frac12$.

The second discovery concerns the zeros. A Brownian path that starts at $0$ returns to $0$ infinitely often in every time interval $[0,\varepsilon]$, yet the set of its zeros
has length zero. It has no isolated points, so it is uncountable, like the Cantor set. In 1955 the British mathematician James Taylor measured its size in the sense of
Hausdorff: dimension exactly $\frac12$. And in 1961 Aryeh Dvoretzky, Paul Erd\H{o}s and Shizuo Kakutani settled a question that sounds absurd, whether a Brownian path
can have a \emph{point of increase}, a time before which it stays below its current value and after which it stays above, for a little while. A continuous function on
$[0,1]$ that is not monotone on any interval may still have such points; the answer for Brownian motion is no.

This chapter proves L\'evy's modulus and the structure of the zero set, with its dimension, using three tools: the zero--one law of Blumenthal for what happens immediately
after time $0$, the strong Markov property of Volume~11 for what happens immediately after a random time, and the reflection principle, now for the joint law of the maximum
and the position.

\begin{reminder}[from Volumes 4, 7 and 11, and Chapter 1]
$B$ is a Brownian motion and $\mathcal F^0_t=\sigma(B_s:s\le t)$. For each $s$, $(B_{s+t}-B_s)_{t\ge0}$ is a Brownian motion independent of $\mathcal F^0_s$, and $c^{-1}B_{c^2t}$ is a Brownian motion
(Volume~4, Proposition~11.5); so is $tB_{1/t}$ (Volume~4, Proposition~11.7). The maximum $M_t=\max_{s\le t}B_s$ has the law of $|B_t|$ (Volume~4, Theorem~11.12). \emph{Strong Markov property}
(Volume~11, Theorem~3.8): if $\tau$ is an almost surely finite stopping time, $(B_{\tau+t}-B_\tau)_{t\ge0}$ is a Brownian motion independent of $\mathcal F_\tau$. Gaussian tails (Volume~11, (3.1)):
$\frac x{1+x^2}\frac{e^{-x^2/2}}{\sqrt{2\pi}}\le P(Z\ge x)\le\frac12e^{-x^2/2}$ for $x>0$. Paths are H\"older of every order $\gamma<\frac12$ (Corollary~\ref{cor:bmholder}). \emph{Hausdorff dimension} (Volume~7, Chapter~12):
if $\mathcal H^s(A)=0$ then $\dim_HA\le s$; and by the \emph{mass distribution principle} (Volume~7, Lemma~12.10), if a finite measure $\mu$ with $\mu(A)>0$ satisfies $\mu(U)\le c(\operatorname{diam}U)^s$ for closed sets
$U$ of small diameter, then $\dim_HA\ge s$.
\end{reminder}

\section{Immediately after time zero, and at the end of time}

Let $\mathcal F^0_{0+}=\bigcap_{t>0}\mathcal F^0_t$: the events that can be decided by watching the path for an arbitrarily short time.

\begin{theorem}[Blumenthal's zero--one law, 1957]\label{thm:blumenthal}
Every event of $\mathcal F^0_{0+}$ has probability $0$ or $1$.
\end{theorem}

\begin{proof}
Let $A\in\mathcal F^0_{0+}$, $0<t_1<\dots<t_k$ and $g:\R^k\to\R$ bounded and continuous. For $0<\varepsilon<t_1$, $A\in\mathcal F^0_\varepsilon$ is independent of the increments after $\varepsilon$, so
\[
\mathbb E\bigl[\mathbf 1_Ag(B_{t_1}-B_\varepsilon,\dots,B_{t_k}-B_\varepsilon)\bigr]=P(A)\,\mathbb E\,g(B_{t_1}-B_\varepsilon,\dots,B_{t_k}-B_\varepsilon).
\]
Let $\varepsilon\to0$: $B_\varepsilon\to0$ and bounded convergence gives
\[\mathbb E[\mathbf 1_Ag(B_{t_1},\dots,B_{t_k})]=P(A)\,\mathbb E g(B_{t_1},\dots,B_{t_k}).\] So $A$ is independent of every vector
$(B_{t_1},\dots,B_{t_k})$ with $t_1>0$, hence (by the $\pi$-$\lambda$ theorem, Volume~11, Lemma~1.1) of $\sigma(B_t:t>0)$, which contains $B_0=\lim B_{1/n}$ and therefore $\mathcal F^0_{0+}$. In particular $A$ is
independent of itself: $P(A)=P(A)^2$.
\end{proof}

\begin{corollary}\label{cor:oscillate}
Almost surely, for every $\varepsilon>0$, $B$ takes positive and negative values in $]0,\varepsilon[$, and hence vanishes there infinitely often.
\end{corollary}

\begin{proof}
Let $\tau_+=\inf\{t>0:B_t>0\}$. The event $\{\tau_+=0\}=\bigcap_n\{\exists q\in\Q\cap]0,\frac1n[:B_q>0\}$ is in $\mathcal F^0_{0+}$, and $P(\tau_+\le\varepsilon)\ge P(B_\varepsilon>0)=\frac12$ for all $\varepsilon$, so
$P(\tau_+=0)\ge\frac12$, hence $=1$. By symmetry the same holds for negative values, and the intermediate value theorem gives zeros between a positive and a negative value.
\end{proof}

Time inversion turns the beginning of time into its end. The \emph{tail} $\sigma$-algebra $\mathcal T=\bigcap_t\sigma(B_s:s\ge t)$ becomes $\mathcal F^0_{0+}$ for the Brownian motion $X_t=tB_{1/t}$, since $B_s=sX_{1/s}$.

\begin{corollary}\label{cor:infinity}
Every event of $\mathcal T$ has probability $0$ or $1$. Almost surely $\sup_tB_t=+\infty$ and $\inf_tB_t=-\infty$; so the path returns to $0$ at arbitrarily large times.
\end{corollary}

\begin{proof}
The first claim is Theorem~\ref{thm:blumenthal} for $X$. Let $S=\sup_tB_t\in[0,\infty]$. By scaling, $cS=\sup_tcB_{t/c^2}$ has the law of $S$ for every $c>0$, so $P(0<S<\infty)=0$: otherwise one of the disjoint intervals $]2^k,2^{k+1}]$, $k\in\Z$, would have positive mass, and by invariance under $x\mapsto2x$ all of them would have that same mass,
adding up to infinity. So $S\in\{0,\infty\}$ almost surely, and the same holds for $S'=\sup_{t\ge1}(B_t-B_1)$, the supremum of the Brownian motion $(B_{1+t}-B_1)_t$, which is independent of $B_1$. On
$\{S=0\}$ we have $B_1\le0$ and $S'\le-B_1<\infty$, so $S'=0$; hence $P(S=0)\le P(B_1\le0)P(S'=0)=\frac12P(S=0)$, and $P(S=0)=0$. Thus $S=+\infty$, and $\inf B=-\infty$ by symmetry. Since
$(B_{n+t}-B_n)_t$ is a Brownian motion, $\sup_{t\ge n}B_t=+\infty$ and $\inf_{t\ge n}B_t=-\infty$ for every $n$, and the path crosses $0$ after time $n$.
\end{proof}

\section{L\'evy's modulus of continuity}

\begin{think}
At a single time, the law of the iterated logarithm bounds $|B_{t+h}-B_t|$ by about $\sqrt{2h\log\log(1/h)}$. Over all $t\in[0,1]$ at once, how much larger can the worst increment be? A
hint: at scale $h$ there are about $1/h$ independent increments, and the largest of $N$ independent standard Gaussians is about $\sqrt{2\log N}$.
\end{think}

\begin{theorem}[L\'evy, 1937]\label{thm:levymodulus}
Almost surely,
\[
\limsup_{h\downarrow0}\ \sup_{0\le t\le1-h}\frac{|B_{t+h}-B_t|}{\sqrt{2h\log(1/h)}}=1 .
\]
\end{theorem}

\begin{proof}
Write $\psi(h)=\sqrt{2h\log(1/h)}$; it is increasing on $]0,e^{-1}]$.

\emph{Lower bound.} Let $0<c<1$, $h=2^{-m}$ and $x_m=c\sqrt{2\log(2^m)}$. The $2^m$ normalized increments $2^{m/2}(B_{k2^{-m}}-B_{(k-1)2^{-m}})$ are independent standard Gaussians, so
\[
P\Bigl(\max_{k\le2^m}|B_{k2^{-m}}-B_{(k-1)2^{-m}}|\le c\,\psi(2^{-m})\Bigr)=\bigl(1-P(|Z|>x_m)\bigr)^{2^m}\le\exp\bigl(-2^mP(|Z|>x_m)\bigr).
\]
By the lower Gaussian tail bound, $P(|Z|>x_m)\ge\frac{c_1}{x_m}e^{-x_m^2/2}=\frac{c_1}{x_m}2^{-c^2m}$ for some $c_1>0$ and $m$ large, so $2^mP(|Z|>x_m)\ge c_1\frac{2^{m(1-c^2)}}{x_m}$, which grows exponentially. The
probabilities are summable, and by Borel--Cantelli, for all large $m$ some increment over an interval of length $2^{-m}$ exceeds $c\psi(2^{-m})$. So the $\limsup$ is at least $c$, for every $c<1$.

\emph{Upper bound.} Let $\varepsilon>0$ and choose $\delta\in]0,1[$ so small that $(1-\delta)(1+\varepsilon)^2>1+\delta$. For $m\ge1$, consider all pairs of dyadic points $j2^{-m}<k2^{-m}$ of $[0,1]$ with
$k-j\le2^{m\delta}$; there are at most $2^{m(1+\delta)}+2^m$ of them. For such a pair, with $h=(k-j)2^{-m}$, the upper Gaussian tail bound gives
\[
P\bigl(|B_{k2^{-m}}-B_{j2^{-m}}|>(1+\varepsilon)\psi(h)\bigr)\le e^{-(1+\varepsilon)^2\log(1/h)}=h^{(1+\varepsilon)^2}\le2^{-m(1-\delta)(1+\varepsilon)^2},
\]
since $h\le2^{-m(1-\delta)}$. The sum over pairs and over $m$ converges, by the choice of $\delta$. By Borel--Cantelli there is a random $m_0$ such that for all $m\ge m_0$ and all such pairs,
\begin{equation}\label{eq:modulusgood}
|B_{k2^{-m}}-B_{j2^{-m}}|\le(1+\varepsilon)\psi\bigl((k-j)2^{-m}\bigr).
\end{equation}
Now let $0\le t<s\le1$ with $h=s-t$ small, and choose $m$ with $2^{-(m+1)(1-\delta)}<h\le2^{-m(1-\delta)}$; for $h$ small, $m\ge m_0$ and $h\ge2\cdot2^{-m}$. Let $t_n=2^{-n}\lceil2^nt\rceil$ and $s_n=2^{-n}\lfloor2^ns\rfloor$
for $n\ge m$. Then $t\le t_m<s_m\le s$, $s_m-t_m\le h$ is a multiple $(k-j)2^{-m}$ with $k-j\le h2^m\le2^{m\delta}$, and~\eqref{eq:modulusgood} with the monotonicity of $\psi$ gives
$|B_{s_m}-B_{t_m}|\le(1+\varepsilon)\psi(h)$. Moreover $t_n-t_{n+1}\in\{0,2^{-n-1}\}$, a pair of neighbours of level $n+1$, and $t_n\to t$; likewise for $s_n$. By~\eqref{eq:modulusgood} with $k-j=1$ and continuity,
\[
|B_t-B_{t_m}|+|B_s-B_{s_m}|\le2\sum_{n>m}(1+\varepsilon)\psi(2^{-n})\le C\sqrt{m2^{-m}},
\]
for a constant $C$, because $\psi(2^{-n})=\sqrt{2n\log2\cdot2^{-n}}$ and the terms decrease geometrically. Since $h>2^{-(m+1)(1-\delta)}$, $\psi(h)^2\ge2^{-(m+1)(1-\delta)}(m+1)(1-\delta)\log2$, and
$\sqrt{m2^{-m}}/\psi(h)\le C'2^{-m\delta/2}\to0$. So for $h$ small enough $|B_s-B_t|\le(1+2\varepsilon)\psi(h)$, and the $\limsup$ is at most $1+2\varepsilon$.
\end{proof}

\begin{figure}[ht]
\centering
\fig{ch02-modulus}
\caption{For one Brownian path sampled at $2^{16}$ points, the largest increment $\max_t|B_{t+h}-B_t|$ over the grid, for $h=2^{-1},\dots,2^{-14}$ (dots), against L\'evy's modulus
$\sqrt{2h\log(1/h)}$ (solid) and against $\sqrt h$ (dashed), on logarithmic scales. The gap between $\sqrt h$ and the dots is the logarithm.}
\label{fig:modulus}
\end{figure}

Compare with the law of the iterated logarithm: at a fixed time the increments are smaller, $\sqrt{2h\log\log(1/h)}$, and by Fubini the times at which L\'evy's modulus is attained form a set of
measure zero (Exercise~\ref{ex:fast}). They do exist, as Steven Orey and James Taylor showed in 1974; they are the \emph{fast points} of the path.

\section{The maximum and the position}

Volume~4 found the law of the maximum $M_t$ by reflecting the path at the first time it reaches a level. With the strong Markov property we can reflect the joint law.

\begin{theorem}[reflection principle, joint form]\label{thm:jointreflection}
Let $a>0$ and $b\le a$. Then $P(M_t\ge a,\ B_t\le b)=P(B_t\ge2a-b)$. Consequently $(M_t,B_t)$ has the density
\[
f(a,b)=\frac{2(2a-b)}{\sqrt{2\pi t^3}}\exp\Bigl(-\frac{(2a-b)^2}{2t}\Bigr)\qquad\text{on }\{a>0,\ b<a\}.
\]
\end{theorem}

\begin{proof}
Let $\tau=\inf\{s:B_s=a\}$, a stopping time, finite by Corollary~\ref{cor:infinity}, and $W_u=B_{\tau+u}-a$, a Brownian motion independent of $\mathcal F_\tau$, so of $\tau$ (Volume~11, Theorem~3.8). On
$\{\tau\le t\}$, $B_t=a+W_{t-\tau}$. The pair $(\tau,W)$ has the law of $(\tau,-W)$, since $W$ and $-W$ are Brownian motions independent of $\tau$. The map $(s,w)\mapsto w(t-s)\mathbf 1_{s\le t}$ is measurable
on $[0,\infty[\times C[0,\infty[$, so
\[
P(\tau\le t,\ W_{t-\tau}\le b-a)=P(\tau\le t,\ -W_{t-\tau}\le b-a)=P(\tau\le t,\ W_{t-\tau}\ge a-b).
\]
The left side is $P(M_t\ge a,B_t\le b)$. On the right, $\{\tau\le t,\ B_t\ge2a-b\}=\{B_t\ge2a-b\}$ because $2a-b\ge a$ forces $\tau\le t$. Differentiating $P(B_t\ge2a-b)=\int_{2a-b}^\infty\frac{e^{-x^2/2t}}{\sqrt{2\pi t}}\dd x$
in $b$ and then in $a$ (with a minus sign) gives the density.
\end{proof}

\begin{corollary}[L\'evy's identity at a fixed time]\label{cor:levyfixed}
For each $t$, $M_t-B_t$ has the law of $|B_t|$, and so has $M_t$. And the time-reversed process $\hat B_u=B_t-B_{t-u}$, $0\le u\le t$, is a Brownian motion on $[0,t]$, with $M_t-B_t=\max_{u\le t}(-\hat B_u)$.
\end{corollary}

\begin{proof}
$\hat B$ is centred Gaussian with continuous paths, $\hat B_0=0$, and for $u\le v$, $\operatorname{Cov}(\hat B_u,\hat B_v)=t-(t-u)-(t-v)+(t-v)=u$; so it is a Brownian motion on $[0,t]$. Then
$M_t-B_t=\max_{s\le t}(B_s-B_t)=\max_{u\le t}(-\hat B_u)$, which has the law of $|B_t|$ by Volume~4, Theorem~11.12 applied to $-\hat B$.
\end{proof}

L\'evy proved much more: the whole \emph{process} $(M_t-B_t)_{t\ge0}$ has the law of $(|B_t|)_{t\ge0}$, and $M_t$ is then the ``local time'' of $|B|$ at zero. We prove this in Chapter~4, with
stochastic calculus, and it will complete the computation of the dimension of the zero set below.

\section{The zero set}

Let $\mathcal Z=\{t\ge0:B_t=0\}$. It is closed, since $B$ is continuous, and it has Lebesgue measure zero almost surely, by Tonelli: $\mathbb E\lambda(\mathcal Z)=\int_0^\infty P(B_t=0)\dd t=0$ (Volume~4, Exercise~11.3).

\begin{theorem}[the zero set]\label{thm:zeroset}
Almost surely, $\mathcal Z$ is a closed, unbounded set of Lebesgue measure zero without isolated points. In particular it is uncountable.
\end{theorem}

\begin{proof}
It is unbounded by Corollary~\ref{cor:infinity}. For rational $q\ge0$ let $\tau_q=\inf\{t\ge q:B_t=0\}$, a stopping time, almost surely finite. By the strong Markov property, $(B_{\tau_q+u})_{u\ge0}$ is a
Brownian motion, since $B_{\tau_q}=0$; by Corollary~\ref{cor:oscillate} it has zeros in every interval $]0,\varepsilon[$. So, outside a null event (a countable union of null events), for every rational $q$ the
zero $\tau_q$ is a limit of zeros from the right.

Now let $z\in\mathcal Z$. If $z=0$, it is a limit of zeros by Corollary~\ref{cor:oscillate}. If $z>0$, either there are zeros in $]z-\varepsilon,z[$ for every $\varepsilon>0$, and $z$ is a limit from the left; or some
interval $]z-\varepsilon,z[$ contains no zero, and then for a rational $q$ in it, $\tau_q=z$, which is a limit from the right. In both cases $z$ is not isolated. A non-empty closed set without isolated points is
uncountable (Volume~4, Exercise~11.3(c)).
\end{proof}

\begin{beauty}
The proof is a small masterpiece of logic. A typical zero is approached from the left, and there is nothing to prove; the zeros that are \emph{not} approached from the left are exactly the
first zeros after rational times, there are only countably many of them, and at each of them the path starts afresh, as if from time $0$. Uncountably many times are controlled by
countably many restarts.
\end{beauty}

\section{Dimension one half}

How big is a set that is uncountable but of length zero? The Cantor set of Volume~7 has dimension $\frac{\log2}{\log3}$. For the zeros of Brownian motion the answer is $\frac12$, and it is
natural: at scale $h$, the path near a zero looks like a path of a random walk with $1/h$ steps, which returns to zero about $\sqrt{1/h}$ times.

\begin{figure}[ht]
\centering
\fig{ch02-zeros}
\caption{The zero set of one Brownian path, seen on $[0,1]$, then on the first $\frac1{16}$, then on the first $\frac1{256}$, each rescaled to the same width; the path is drawn above each strip, rescaled to the same height, and the rust ticks mark the zeros (clusters closer than the
resolution are merged). By Brownian scaling the three strips have the same law: the zero set is statistically self-similar, a random fractal.}
\label{fig:zeros}
\end{figure}

We also consider the \emph{record set} $\mathcal R=\{t\ge0:B_t=M_t\}$, the times at which the path is at its running maximum. It is closed, and it is where the continuous non-decreasing function $M$
increases: on each component of the open set $\{t:B_t<M_t\}$, $M$ is constant.

\begin{lemma}\label{lem:hitprob}
For $0<a$ and $h>0$,
\[
P\bigl(\mathcal Z\cap[a,a+h]\neq\varnothing\bigr)\le\frac4\pi\sqrt{\frac ha}\qquad\text{and}\qquad P\bigl(\mathcal R\cap[a,a+h]\neq\varnothing\bigr)\le\frac2\pi\sqrt{\frac ha}.
\]
\end{lemma}

\begin{proof}
Let $D^+=\max_{s\in[a,a+h]}(B_s-B_a)$ and $D^-=\max_{s\in[a,a+h]}(B_a-B_s)$. They are independent of $\mathcal F^0_a$ and each has the law of $|B_h|$, with $\mathbb E|B_h|=\sqrt{2h/\pi}$. If $B$ vanishes in $[a,a+h]$,
then $|B_a|\le D^++D^-$. Since $B_a$ has density at most $\frac1{\sqrt{2\pi a}}$, $P(|B_a|\le x)\le\frac{2x}{\sqrt{2\pi a}}$, and conditioning on $D^\pm$ (independent of $B_a$),
\[
P(\mathcal Z\cap[a,a+h]\neq\varnothing)\le\mathbb E\frac{2(D^++D^-)}{\sqrt{2\pi a}}=\frac{4}{\sqrt{2\pi a}}\sqrt{\frac{2h}\pi}=\frac4\pi\sqrt{\frac ha}.
\]
For the record set, $\mathcal R$ meets $[a,a+h]$ only if $\max_{[a,a+h]}B\ge M_a$, that is $D^+\ge M_a-B_a$. By Corollary~\ref{cor:levyfixed}, $M_a-B_a$ has the law of $|B_a|$, and it is $\mathcal F^0_a$-measurable,
hence independent of $D^+$. The same computation with $D^+$ alone gives $\frac2\pi\sqrt{h/a}$.
\end{proof}

\begin{theorem}[Taylor, 1955]\label{thm:dimension}
Almost surely, $\dim_H(\mathcal Z\cap[0,1])\le\frac12$ and $\dim_H(\mathcal R\cap[0,1])=\frac12$.
\end{theorem}

\begin{proof}
\emph{Upper bounds.} Let $A$ be $\mathcal Z$ or $\mathcal R$. Cut $[0,1]$ into $n$ intervals $I_k=[\frac{k-1}n,\frac kn]$ and let $N_n$ be the number of them that meet $A$. By Lemma~\ref{lem:hitprob} with $a=\frac{k-1}n$ and
$h=\frac1n$,
\[
\mathbb E N_n\le1+\frac4\pi\sum_{j=1}^{n-1}\frac1{\sqrt j}\le1+\frac8\pi\sqrt n .
\]
Let $s>\frac12$. Covering $A\cap[0,1]$ by the $N_n$ intervals of length $\frac1n$ gives $\mathcal H^s_{1/n}(A\cap[0,1])\le N_nn^{-s}$, so for every $\delta>0$, $\mathcal H^s_\delta(A\cap[0,1])\le\liminf_nN_nn^{-s}$. By Fatou,
$\mathbb E\liminf N_nn^{-s}\le\liminf\mathbb E N_nn^{-s}=0$. Hence $\mathcal H^s_\delta(A\cap[0,1])=0$ for all $\delta$, almost surely, so $\mathcal H^s(A\cap[0,1])=0$ and $\dim_H(A\cap[0,1])\le s$. Let $s\downarrow\frac12$ along a sequence.

\emph{Lower bound for $\mathcal R$.} Let $\mu$ be the Stieltjes measure of $M$ on $[0,1]$: $\mu([u,v])=M_v-M_u$. It is carried by $\mathcal R$, since $M$ is constant on the components of the complement,
and $\mu(\mathcal R\cap[0,1])=M_1>0$ almost surely. Fix $\gamma<\frac12$. By Corollary~\ref{cor:bmholder} there is a random $K_\gamma$ with $|B_v-B_u|\le K_\gamma|v-u|^\gamma$ on $[0,1]$. If $U\subseteq[0,1]$ is closed with diameter $h$,
it lies in an interval $[u,u+h]$, and
\[
\mu(U)\le M_{u+h}-M_u\le\max_{s\in[u,u+h]}(B_s-B_u)\le K_\gamma h^\gamma,
\]
because $M_{u+h}$ is either $M_u$ or a value $B_s$ with $s\in[u,u+h]$, and $M_u\ge B_u$. By the mass distribution principle, $\dim_H(\mathcal R\cap[0,1])\ge\gamma$, for every $\gamma<\frac12$.
\end{proof}

The lower bound for $\mathcal Z$ needs a natural measure on the zero set, the analogue of $\mathrm dM$ for $\mathcal R$. L\'evy's theorem of Chapter~4, that $M-B$ and $|B|$ have the same law as processes, will
carry the record set onto the zero set and give $\dim_H(\mathcal Z\cap[0,1])=\frac12$ (Corollary~\ref{cor:zerodim}). The measure that does the job is the local time at $0$.

\begin{remark}
Taylor (1955) proved more: the zero set has $\sigma$-finite but not finite $\mathcal H^{1/2}$-measure, and later work found the exact gauge function $\sqrt{h\log\log(1/h)}$ for which its Hausdorff measure is
a constant multiple of the local time. The graph $\{(t,B_t):t\le1\}$ has dimension $\frac32$ (Exercise~\ref{ex:graph} for the upper bound) and the image of a set $E\subseteq[0,\infty[$ under planar Brownian motion has
dimension $2\dim_HE$, almost surely for all $E$ at once (Kaufman, 1969). Finally, Dvoretzky, Erd\H{o}s and Kakutani (1961) proved that almost surely Brownian motion has no point of increase; the shortest proof
known, by Yuval Peres, goes through random walks and is in M\"orters and Peres, Chapter~5.
\end{remark}

\section{Exercises}

\begin{exercise}[$\star$]\label{ex:limsupsqrt}
Show that almost surely $\limsup_{t\downarrow0}B_t/\sqrt t=+\infty$ and $\limsup_{t\to\infty}B_t/\sqrt t=+\infty$. (For each $c$, consider $\{B_t>c\sqrt t\text{ for arbitrarily small }t\}$ and use
Theorem~\ref{thm:blumenthal}.) Compare with the law of the iterated logarithm.
\end{exercise}

\begin{exercise}[$\star\star$]\label{ex:localmax}
(a) Show that for rational $0\le p<q<r<s$, $\max_{[p,q]}B\neq\max_{[r,s]}B$ almost surely. (b) Deduce that almost surely $B$ attains its maximum on $[0,1]$ at a single point, that the values at local maxima
are all distinct, and that the set of times of local maxima is countable and dense. (c) Show that $B$ is monotone on no interval.
\end{exercise}

\begin{exercise}[$\star\star$]\label{ex:bridgemax}
(a) From Theorem~\ref{thm:jointreflection}, show that conditionally on $B_t=b$, the probability that $M_t\ge a$ is $\exp\bigl(-\frac{2a(a-b)}t\bigr)$ for $a\ge\max(b,0)$. (b) Deduce that the maximum of the Brownian bridge
on $[0,1]$ has the distribution function $1-e^{-2a^2}$. (c) Check from the joint density that $M_t-B_t$ and $M_t$ each have the law of $|B_t|$, but are not independent.
\end{exercise}

\begin{exercise}[$\star\star$]\label{ex:nozero}
Show that for $0<a<b$, $P(\mathcal Z\cap]a,b[=\varnothing)=\frac2\pi\arcsin\sqrt{a/b}$. (Condition on $B_a$ and use the reflection principle for the Brownian motion after $a$; then compute a Gaussian probability in polar
coordinates.)
\end{exercise}

\begin{exercise}[$\star\star$]\label{ex:graph}
Show that almost surely the graph $G=\{(t,B_t):0\le t\le1\}$ has $\dim_HG\le\frac32$: cover it, for each $\gamma<\frac12$, by squares of side $\frac1n$ stacked above each interval $I_k$, and count them with the H\"older bound.
\end{exercise}

\begin{exercise}[$\star\star$]\label{ex:walkzeros}
Let $S$ be the simple random walk and $Z_n=\#\{k\le n:S_k=0\}$. Show that $\mathbb E Z_{2n}=\sum_{k\le n}\binom{2k}k4^{-k}\sim2\sqrt{n/\pi}$. Interpret: the zeros of the walk among $N$ steps are about $\sqrt N$ in number,
as a set of dimension $\frac12$ should be.
\end{exercise}

\begin{exercise}[$\star\star\star$]\label{ex:fast}
Call $t\in[0,1[$ a \emph{fast point} if $\limsup_{h\downarrow0}|B_{t+h}-B_t|/\sqrt{2h\log(1/h)}>0$. Show that almost surely the set of fast points has Lebesgue measure zero. (Use Fubini and the law of the iterated
logarithm of Volume~11, Theorem~3.18, at each fixed $t$, after checking that the set of fast points is measurable in $(t,\omega)$.)
\end{exercise}

\section{Hints}

\begin{hint}{ex:limsupsqrt}
$P(B_t>c\sqrt t)=P(Z>c)>0$ for every $t$; the event that this happens for $t$ arbitrarily close to $0$ has probability at least $P(Z>c)$ by Fatou applied to indicators, and lies in $\mathcal F^0_{0+}$.
\end{hint}

\begin{hint}{ex:localmax}
$\max_{[r,s]}B-\max_{[p,q]}B=(B_r-B_q)+\bigl(\max_{[r,s]}B-B_r\bigr)-\bigl(\max_{[p,q]}B-B_q\bigr)$, and $B_r-B_q$ is Gaussian independent of the rest.
\end{hint}

\begin{hint}{ex:bridgemax}
Divide the joint density by the density of $B_t$, and integrate in $a$.
\end{hint}

\begin{hint}{ex:nozero}
$P(\text{no zero in }]a,b[\mid B_a=x)=P(|B_{b-a}|<|x|)$.
\end{hint}

\begin{hint}{ex:graph}
Above $I_k$, the graph fits in a column of height at most $K_\gamma n^{-\gamma}$, covered by at most $K_\gamma n^{1-\gamma}+2$ squares of side $\frac1n$.
\end{hint}

\begin{hint}{ex:walkzeros}
$\binom{2k}k4^{-k}\sim\frac1{\sqrt{\pi k}}$ (Stirling), and $\sum_{k\le n}k^{-1/2}\sim2\sqrt n$.
\end{hint}

\begin{hint}{ex:fast}
$F=\{(t,\omega):\limsup_{h\downarrow0,\,h\in\Q}|B_{t+h}-B_t|/\psi(h)>0\}$ is measurable, and by continuity the $\limsup$ over rational $h$ equals the one over all $h$. For fixed $t$, $\sqrt{2h\log\log(1/h)}/\psi(h)\to0$.
\end{hint}

\section{Solutions}

\begin{solution}{ex:limsupsqrt}
\idea The ratio $B_t/\sqrt t$ is standard Gaussian at every time, so exceeding a fixed level $c$ has a fixed positive probability at every scale. An event decided by arbitrarily small times has
probability $0$ or $1$, so positive means $1$.
\writeup Fix $c>0$ and let
\[A_c=\bigcap_n\bigcup_{t\in\Q\cap]0,1/n[}\{B_t>c\sqrt t\},\]
which lies in $\mathcal F^0_{0+}$. The events $E_n=\{B_{1/n}>c/\sqrt n\}$ have probability $P(Z>c)$, and
$A_c\supseteq\limsup E_n$, so $P(A_c)\ge P(\limsup E_n)\ge\limsup P(E_n)=P(Z>c)>0$ (Fatou for indicators). By Theorem~\ref{thm:blumenthal}, $P(A_c)=1$. Intersecting over $c=1,2,\dots$ gives
$\limsup_{t\downarrow0}B_t/\sqrt t=+\infty$. At infinity, apply this to $X_t=tB_{1/t}$: $X_t/\sqrt t=\sqrt tB_{1/t}=B_s/\sqrt s$ with $s=1/t$. Both statements are weaker than the law of the iterated logarithm, which gives the
exact rate $\sqrt{2t\log\log(1/t)}$ (at $0$, via time inversion), but they need only the zero--one law.
\end{solution}

\begin{solution}{ex:localmax}
\idea Two maxima over disjoint intervals differ by a Gaussian increment plus independent terms, and a Gaussian plus an independent variable has no atoms. Countably many pairs of rational intervals see every pair
of local maxima.
\writeup (a) In the hint's decomposition, $X=B_r-B_q$ is $N(0,r-q)$ and independent of $Y=(\max_{[r,s]}B-B_r)-(\max_{[p,q]}B-B_q)$, which depends on the path before $q$ and on the increments after $r$. So
$P(X+Y=0)=\mathbb E\,P(X=-Y\mid Y)=0$. (b) If the maximum on $[0,1]$ were attained at $t_1<t_2$, pick rationals $p<q<r<s$ with $t_1\in[p,q]$, $t_2\in[r,s]$, inside $[0,1]$ (if $t_1=0$ use $p=0$): then
$\max_{[p,q]}B=\max_{[r,s]}B$, which by (a) and countability has probability zero. A local maximum at $t$ is the maximum on some rational interval $[p,q]\ni t$, and two different local maxima with the same value would
again give equal maxima over disjoint rational intervals, so the values are distinct. Each rational interval has a single maximum point (the argument of (b) on $[p,q]$), and every local maximum time is the maximum
point of some rational interval, so there are countably many. Density: every interval $[p,q]$ contains the point where $B$ attains its maximum over it, which is a local maximum if it is interior; it is interior
almost surely, because by Corollary~\ref{cor:oscillate} applied at $p$ and (by time reversal) at $q$, the path exceeds $B_p$ just after $p$ and exceeds $B_q$ just before $q$. (c) A function monotone on an interval has no local
maximum in its interior, but local maxima are dense.
\end{solution}

\begin{solution}{ex:bridgemax}
\idea Conditioning on the endpoint turns the joint density into a ratio of Gaussians, and the ratio simplifies to one exponential: this is the famous formula behind the Kolmogorov--Smirnov test.
\writeup (a) For $a>\max(b,0)$, $P(M_t\ge a,B_t\in\dd b)=\frac{1}{\sqrt{2\pi t}}e^{-(2a-b)^2/2t}\dd b$ (differentiate $P(B_t\ge2a-b)$ in $b$). Dividing by $P(B_t\in\dd b)=\frac1{\sqrt{2\pi t}}e^{-b^2/2t}\dd b$ gives
$\exp(-\frac{(2a-b)^2-b^2}{2t})=\exp(-\frac{2a(a-b)}t)$. (b) The Brownian bridge is Brownian motion on $[0,1]$ conditioned on $B_1=0$ (its law is that of $B$ given $B_1=0$, since $B_t-tB_1$ is independent of $B_1$ and equals $B$ when $B_1=0$); with $t=1$, $b=0$: $P(\max\ge a)=e^{-2a^2}$. (c) The law of $M_t-B_t$: substitute $y=a-b>0$ and integrate the density over $a>0$ at fixed $y$: $\int_0^\infty\frac{2(a+y)}{\sqrt{2\pi t^3}}e^{-(a+y)^2/2t}\dd a=\frac2{\sqrt{2\pi t}}e^{-y^2/2t}$, the density of $|B_t|$. Not independent: in the variables $(a,y)=(M_t,M_t-B_t)$ the joint density is $\frac{2(a+y)}{\sqrt{2\pi t^3}}e^{-(a+y)^2/2t}$, which tends to $0$ at $(0,0)$, while the product of the two marginal densities is
$\frac2{\pi t}e^{-(a^2+y^2)/2t}$, equal to $\frac2{\pi t}$ there. Near $(0,0)$ the pair is less likely than independence would allow: a path whose maximum is tiny cannot have fallen far below it.
\end{solution}

\begin{solution}{ex:nozero}
\idea After time $a$, the path starting at $x$ avoids zero until $b$ exactly when the Brownian motion after $a$ does not drop by $|x|$ in time $b-a$, and the reflection principle gives that probability. The final
integral is the probability that a planar Gaussian vector lies in a cone.
\writeup By the Markov property at $a$, given $B_a=x$ the path after $a$ is $x+W$ with $W$ a Brownian motion independent of $B_a$. It has no zero in $]a,b[$ if and only if $\max_{u\le b-a}(-\mathrm{sgn}(x)W_u)<|x|$, which by
Volume~4, Theorem~11.12 has probability $P(|W_{b-a}|<|x|)$. So $P(\text{no zero})=P(|\sqrt{b-a}\,Z'|<|\sqrt a\,Z|)$ with $Z,Z'$ independent standard Gaussians. The vector $(Z,Z')$ has a rotation invariant law, and the event
$|Z'|<\sqrt{a/(b-a)}\,|Z|$ is a double cone of half-angle $\theta=\arctan\sqrt{a/(b-a)}$ around the horizontal axis; its probability is $\frac{4\theta}{2\pi}=\frac2\pi\theta$. Finally $\tan\theta=\sqrt{a/(b-a)}$ means $\sin\theta=\sqrt{a/b}$.
\end{solution}

\begin{solution}{ex:graph}
\idea Hölder continuity says the graph over a short interval fits in a thin column; count small squares and compare with $n^s$.
\writeup Fix $\gamma<\frac12$. Above $I_k$ the graph lies in $I_k\times[m_k,m_k+K_\gamma n^{-\gamma}]$, where $m_k$ is the minimum there; this rectangle is covered by at most $K_\gamma n^{1-\gamma}+1$ squares of side
$\frac1n$, of diameter $\frac{\sqrt2}n$. So $\mathcal H^s_{\sqrt2/n}(G)\le n(K_\gamma n^{1-\gamma}+1)(\sqrt2/n)^s$, which tends to $0$ when $s>2-\gamma$. Hence $\dim_HG\le2-\gamma$ for every $\gamma<\frac12$, so $\le\frac32$.
\end{solution}

\begin{solution}{ex:walkzeros}
\idea Linearity of expectation turns a count into a sum of probabilities, and each probability is the central binomial coefficient of Volume~4.
\writeup $\mathbb E Z_{2n}=\sum_{k=1}^{2n}P(S_k=0)=\sum_{j=1}^nu_{2j}$ with $u_{2j}=\binom{2j}j4^{-j}\sim(\pi j)^{-1/2}$. A sum of positive terms equivalent to $(\pi j)^{-1/2}$, with a divergent sum, is equivalent to
$\sum_{j\le n}(\pi j)^{-1/2}\sim2\sqrt{n/\pi}$. So among $N=2n$ steps there are about $\sqrt{2N/\pi}$ zeros. Seen at scale $h=\frac1N$, a set of dimension $d$ meets about $h^{-d}$ intervals: here $h^{-1/2}$, as
Theorem~\ref{thm:dimension} predicts for the limit.
\end{solution}

\begin{solution}{ex:fast}
\idea For each fixed time the law of the iterated logarithm makes the ratio tend to zero, because $\log\log$ is much smaller than $\log$. Fubini turns ``for each $t$, almost surely'' into ``almost surely, for almost every $t$''.
\writeup Let $\psi(h)=\sqrt{2h\log(1/h)}$. For each rational $h$, $(t,\omega)\mapsto|B_{t+h}-B_t|/\psi(h)$ is measurable (Volume~4, Exercise~11.3(a), for the joint measurability of $(t,\omega)\mapsto B_t$), so
$L(t,\omega)=\limsup_{h\downarrow0,h\in\Q}|B_{t+h}-B_t|/\psi(h)$ is measurable, and by continuity of $h\mapsto B_{t+h}$ it equals the $\limsup$ over all $h$. For fixed $t$, $(B_{t+h}-B_t)_h$ is a Brownian motion,
and by time inversion and Volume~11, Theorem~3.18, $|B_{t+h}-B_t|\le2\sqrt{h\log\log(1/h)}$ for all small $h$, almost surely; since $\sqrt{h\log\log(1/h)}/\psi(h)\to0$, $L(t,\cdot)=0$ almost surely. By Tonelli,
$\mathbb E\,\lambda\{t:L(t)>0\}=\int_0^1P(L(t)>0)\dd t=0$.
\end{solution}

\section*{Sources and further reading}

\begin{itemize}
\item \emph{History.} P.~L\'evy, \emph{Th\'eorie de l'addition des variables al\'eatoires}, Gauthier-Villars, 1937 (the modulus of continuity), and \emph{Processus stochastiques et mouvement brownien},
Gauthier-Villars, 1948. M.~Lo\`eve, ``Paul L\'evy, 1886--1971'', \emph{Annals of Probability} 1 (1973), 1--18, the obituary paraphrased in the opening. R.~M.~Blumenthal, ``An extended Markov property'',
\emph{Transactions of the American Mathematical Society} 85 (1957), 52--72. S.~J.~Taylor, ``The $\alpha$-dimensional measure of the graph and set of zeros of a Brownian path'', \emph{Proceedings of the Cambridge
Philosophical Society} 51 (1955), 265--274. A.~Dvoretzky, P.~Erd\H{o}s and S.~Kakutani, ``Nonincrease everywhere of the Brownian motion process'', \emph{Proceedings of the Fourth Berkeley Symposium} 2 (1961),
103--116. S.~Orey and S.~J.~Taylor, ``How often on a Brownian path does the law of iterated logarithm fail?'', \emph{Proceedings of the London Mathematical Society} 28 (1974), 174--192. R.~Kaufman, ``Une propri\'et\'e
m\'etrique du mouvement brownien'', \emph{Comptes Rendus de l'Acad\'emie des Sciences de Paris} 268 (1969), 727--728.
\item P.~M\"orters and Y.~Peres, \emph{Brownian Motion}, Cambridge University Press, 2010, Chapters~1, 2, 4 and~5. L\'evy's modulus, the zero set, Hausdorff dimension and points of increase; the route followed here.
\item K.~It\^o and H.~P.~McKean, \emph{Diffusion Processes and their Sample Paths}, Springer, 1965. The classic account of L\'evy's discoveries, made rigorous.
\item D.~Revuz and M.~Yor, \emph{Continuous Martingales and Brownian Motion}, 3rd ed., Springer, 1999, Chapters~I--III. Path properties in full generality.
\item K.~Falconer, \emph{Fractal Geometry}, 3rd ed., Wiley, 2014, Chapter~16. Dimensions of Brownian sets, with the energy method.
\end{itemize}
